% Options for packages loaded elsewhere
\PassOptionsToPackage{unicode}{hyperref}
\PassOptionsToPackage{hyphens}{url}
\PassOptionsToPackage{space}{xeCJK}
\documentclass[
  a4paper,
]{article}
\usepackage{xcolor}
\usepackage[margin=22mm]{geometry}
\usepackage{amsmath,amssymb}
\setcounter{secnumdepth}{-\maxdimen} % remove section numbering
\usepackage{iftex}
\ifPDFTeX
  \usepackage[T1]{fontenc}
  \usepackage[utf8]{inputenc}
  \usepackage{textcomp} % provide euro and other symbols
\else % if luatex or xetex
  \usepackage{unicode-math} % this also loads fontspec
  \defaultfontfeatures{Scale=MatchLowercase}
  \defaultfontfeatures[\rmfamily]{Ligatures=TeX,Scale=1}
\fi
\usepackage{lmodern}
\ifPDFTeX\else
  % xetex/luatex font selection
  \setmainfont[]{Times New Roman}
  \ifXeTeX
    \usepackage{xeCJK}
    \setCJKmainfont[]{SimSun}
  \fi
  \ifLuaTeX
    \usepackage[]{luatexja-fontspec}
    \setmainjfont[]{SimSun}
  \fi
\fi
% Use upquote if available, for straight quotes in verbatim environments
\IfFileExists{upquote.sty}{\usepackage{upquote}}{}
\IfFileExists{microtype.sty}{% use microtype if available
  \usepackage[]{microtype}
  \UseMicrotypeSet[protrusion]{basicmath} % disable protrusion for tt fonts
}{}
\makeatletter
\@ifundefined{KOMAClassName}{% if non-KOMA class
  \IfFileExists{parskip.sty}{%
    \usepackage{parskip}
  }{% else
    \setlength{\parindent}{0pt}
    \setlength{\parskip}{6pt plus 2pt minus 1pt}}
}{% if KOMA class
  \KOMAoptions{parskip=half}}
\makeatother
\usepackage{color}
\usepackage{fancyvrb}
\newcommand{\VerbBar}{|}
\newcommand{\VERB}{\Verb[commandchars=\\\{\}]}
\DefineVerbatimEnvironment{Highlighting}{Verbatim}{commandchars=\\\{\}}
% Add ',fontsize=\small' for more characters per line
\newenvironment{Shaded}{}{}
\newcommand{\AlertTok}[1]{\textcolor[rgb]{1.00,0.00,0.00}{\textbf{#1}}}
\newcommand{\AnnotationTok}[1]{\textcolor[rgb]{0.38,0.63,0.69}{\textbf{\textit{#1}}}}
\newcommand{\AttributeTok}[1]{\textcolor[rgb]{0.49,0.56,0.16}{#1}}
\newcommand{\BaseNTok}[1]{\textcolor[rgb]{0.25,0.63,0.44}{#1}}
\newcommand{\BuiltInTok}[1]{\textcolor[rgb]{0.00,0.50,0.00}{#1}}
\newcommand{\CharTok}[1]{\textcolor[rgb]{0.25,0.44,0.63}{#1}}
\newcommand{\CommentTok}[1]{\textcolor[rgb]{0.38,0.63,0.69}{\textit{#1}}}
\newcommand{\CommentVarTok}[1]{\textcolor[rgb]{0.38,0.63,0.69}{\textbf{\textit{#1}}}}
\newcommand{\ConstantTok}[1]{\textcolor[rgb]{0.53,0.00,0.00}{#1}}
\newcommand{\ControlFlowTok}[1]{\textcolor[rgb]{0.00,0.44,0.13}{\textbf{#1}}}
\newcommand{\DataTypeTok}[1]{\textcolor[rgb]{0.56,0.13,0.00}{#1}}
\newcommand{\DecValTok}[1]{\textcolor[rgb]{0.25,0.63,0.44}{#1}}
\newcommand{\DocumentationTok}[1]{\textcolor[rgb]{0.73,0.13,0.13}{\textit{#1}}}
\newcommand{\ErrorTok}[1]{\textcolor[rgb]{1.00,0.00,0.00}{\textbf{#1}}}
\newcommand{\ExtensionTok}[1]{#1}
\newcommand{\FloatTok}[1]{\textcolor[rgb]{0.25,0.63,0.44}{#1}}
\newcommand{\FunctionTok}[1]{\textcolor[rgb]{0.02,0.16,0.49}{#1}}
\newcommand{\ImportTok}[1]{\textcolor[rgb]{0.00,0.50,0.00}{\textbf{#1}}}
\newcommand{\InformationTok}[1]{\textcolor[rgb]{0.38,0.63,0.69}{\textbf{\textit{#1}}}}
\newcommand{\KeywordTok}[1]{\textcolor[rgb]{0.00,0.44,0.13}{\textbf{#1}}}
\newcommand{\NormalTok}[1]{#1}
\newcommand{\OperatorTok}[1]{\textcolor[rgb]{0.40,0.40,0.40}{#1}}
\newcommand{\OtherTok}[1]{\textcolor[rgb]{0.00,0.44,0.13}{#1}}
\newcommand{\PreprocessorTok}[1]{\textcolor[rgb]{0.74,0.48,0.00}{#1}}
\newcommand{\RegionMarkerTok}[1]{#1}
\newcommand{\SpecialCharTok}[1]{\textcolor[rgb]{0.25,0.44,0.63}{#1}}
\newcommand{\SpecialStringTok}[1]{\textcolor[rgb]{0.73,0.40,0.53}{#1}}
\newcommand{\StringTok}[1]{\textcolor[rgb]{0.25,0.44,0.63}{#1}}
\newcommand{\VariableTok}[1]{\textcolor[rgb]{0.10,0.09,0.49}{#1}}
\newcommand{\VerbatimStringTok}[1]{\textcolor[rgb]{0.25,0.44,0.63}{#1}}
\newcommand{\WarningTok}[1]{\textcolor[rgb]{0.38,0.63,0.69}{\textbf{\textit{#1}}}}
\usepackage{longtable,booktabs,array}
\newcounter{none} % for unnumbered tables
\usepackage{calc} % for calculating minipage widths
% Correct order of tables after \paragraph or \subparagraph
\usepackage{etoolbox}
\makeatletter
\patchcmd\longtable{\par}{\if@noskipsec\mbox{}\fi\par}{}{}
\makeatother
% Allow footnotes in longtable head/foot
\IfFileExists{footnotehyper.sty}{\usepackage{footnotehyper}}{\usepackage{footnote}}
\makesavenoteenv{longtable}
\setlength{\emergencystretch}{3em} % prevent overfull lines
\providecommand{\tightlist}{%
  \setlength{\itemsep}{0pt}\setlength{\parskip}{0pt}}
\usepackage{bookmark}
\IfFileExists{xurl.sty}{\usepackage{xurl}}{} % add URL line breaks if available
\urlstyle{same}
\hypersetup{
  hidelinks,
  pdfcreator={LaTeX via pandoc}}

\author{GPT}
\date{}

\begin{document}

\section{四列重力四子棋的统一回应算法与交替段压缩}\label{ux56dbux5217ux91cdux529bux56dbux5b50ux68cbux7684ux7edfux4e00ux56deux5e94ux7b97ux6cd5ux4e0eux4ea4ux66ffux6bb5ux538bux7f29}

GPT

2026年11月8日

\subsubsection{摘要}\label{ux6458ux8981}

本文将四列重力四子棋的覆盖策略组织为双方共用、适用于全部有限高度的确定回应规则。该规则只检查当前完整棋盘、自己的颜色和对方上一手列，不保存阶段、已选目标或主动费用，不展开后续博弈树。证明分离当前局面的几何证书、同列覆盖闭包、满列出口和无证书开局阶段，精确规定优先顺序。可执行的
Lean 4
回应函数已完成双方不败性和主动预算最优值的参数化证明，历史前缀语义也已接通。进一步的纸面计数证明把所有合法候选棋盘的最大交替段总数收紧为九；一个有序区间扫描把单个制胜目标搜索从段组合枚举改为线性扫描。完整压缩引擎与已验收函数的精化、以及成本定理尚未形式化，不能从有限对照测试推得。

关键词：策略证书；局面不变量；交替段；区间扫描；Lean 4。

\subsubsection{与数学论文和既有工作的关系}\label{ux4e0eux6570ux5b66ux8bbaux6587ux548cux65e2ux6709ux5de5ux4f5cux7684ux5173ux7cfb}

配套数学论文 {[}4{]}
证明一至四列的空盘和棋及主动换列次数最优值。本文展开四列有限棋盘上的统一策略执行与压缩表示。Allis
{[}1{]} 的配对、Claimeven
和偶数容量释放是相关背景；当前规则根据局面检查选择回应，不搜索后续博弈树。无限高度不属于本算法的输入，不能把
\(H=\infty\) 代入余量奇偶；相关研究见 {[}2{]}。可执行定义与证明见仓库
{[}5{]}，区间扫描及有限检查见 {[}6{]}。

\subsection{1 游戏
计数与算法输入}\label{ux6e38ux620f-ux8ba1ux6570ux4e0eux7b97ux6cd5ux8f93ux5165}

\subsubsection{1.1
首次胜与合法策略}\label{ux9996ux6b21ux80dcux4e0eux5408ux6cd5ux7b56ux7565}

棋盘固定四列，高度为有限非负整数\(H\)。列从左到右编号1至4，行从底到顶编号1至\(H\)。每次在未满列的最低空格落子。横、竖和两种斜向连续四格都计胜；任一方第一次形成四连，游戏立即结束，败方没有补救回应。无胜满盘为和棋，\(H=0\)直接是满盘和棋。

双方始终称先手、后手。列字符串从底至顶写，``先''\,``后''分别缩写其棋子，\(\varnothing\)为空列。下文没有因高度改变双方名称。越界的线或目标一律不计。

固定策略必须覆盖任意对手续法，而不只是两份算法互相对弈。完整不败证明需先排除对方刚落那一手的首胜，再证明自己的回应合法、安全，不能倒过来。

\subsubsection{1.2 换列次数}\label{ux6362ux5217ux6b21ux6570}

对某一方自己的回应计数：若对方刚落的列尚未满，却去另一列回应，称主动换列；若对方刚填满原列后才去别列，称被动换列。同列覆盖不计换列。先手首子不是回应，不计入。自己立即获胜的那枚回应仍按落子前原列是否满来计数，然后停止。

最坏情形主动预算是固定不败策略对所有合法对手续法所需主动次数的上确界；在全部不败策略中取最小值，才是最优主动预算。不能仅击败某一种满列换列选择就证明这一最优值。

\subsubsection{1.3
当前局面输入}\label{ux5f53ux524dux5c40ux9762ux8f93ux5165}

先手在\(H>0\)时首下\((1,1)\)。后手从空盘等待先手第一手。此后双方的回应函数相同，输入为：

\begin{itemize}
\tightlist
\item
  \(H\)与当前完整彩色棋盘，当前是先手还是后手回应；
\item
  对方上一手的位置，记其列为\(c\)、行号为\(r\)。
\end{itemize}

当前棋盘是对方刚落之后的棋盘。重力下\(r\)也可由上一手列的当前高度读出，但必须知道上一手是哪一列；本文不声称只看棋子分布而不读上一手信息。

``局部''在本文指只检查当前棋盘及候选的一手落子，不展开后续博弈树，不是指只看上一手附近几格。所有列高、保护格、例外集合和证书都从当前棋盘重新计算；不保存等待阶段、已用额度、之前选中的目标或之前的证书类别。

以下证书都检查``自己刚回应以后、下一手轮到另一方''的非终局局面。检查先手证书时，自己的格就是先手格，另一方的格就是后手格；检查后手证书时反之。只改变待检查的棋子集合，不交换先手、后手的名称。

\subsection{2
支撑例外与共同安全前提}\label{ux652fux6491ux4f8bux5916ux4e0eux5171ux540cux5b89ux5168ux524dux63d0}

\subsubsection{2.1 支撑下移}\label{ux652fux6491ux4e0bux79fb}

\textbf{引理1。}
若另一方的一条非底、非竖直四连，每枚棋正下方都是本次检查一方的棋，则另一方不能由此首先获胜。

\textbf{证明。}
将整条线下移一行，得到本次检查一方的四连。按重力，这四枚支撑棋在另一方最后落子以前都已存在。因此早已有胜者，与另一方现在首胜矛盾。\(\square\)

\subsubsection{2.2
完整例外与未解决线}\label{ux5b8cux6574ux4f8bux5916ux4e0eux672aux89e3ux51b3ux7ebf}

对一个非终局候选局面，从另一方的棋子完整登记例外集合\(E\)：

\begin{enumerate}
\def\labelenumi{\arabic{enumi}.}
\tightlist
\item
  每个非底的另一方棋，如果其正下方也是另一方棋，登记上格。
\item
  每个非满、由另一方封顶的列，登记该列下一空格。
\end{enumerate}

第二类是下一次进入该列可能产生的支撑例外。满列上方不在棋盘内，不能登记为潜在例外。还要求每个非满列顶端连续另一方棋串至多两枚；已有任何一方四连或已经满盘时，应按终局处理，不再用非终局证书代替终局判断。

经过底行或\(E\)、而尚未含本次检查一方棋子的非竖直四连，组成未解决线集合\(U\)。其他非竖直线已有固定棋子阻断，或由引理1排除另一方首胜。底横线及两条碰到底行的斜线简称底部线；\(H\le 3\)时斜线不存在。

\subsubsection{2.3
覆盖下的例外变化}\label{ux8986ux76d6ux4e0bux7684ux4f8bux5916ux53d8ux5316}

当另一方进入未满列而本方随后同列盖，下一轮的例外不会超出原\(E\)；原来潜在的例外可能变成实际例外，不能漏掉。另一方填满列后，本方在别列加子，也不会创造新的另一方例外。初始非满另一方顶串至多二，最多再增加一枚，随后被本方或顶边封住，故不能新成长为竖四。

这些事实在下一轮仍可重新计算。已阻断的线不会重新打开，新的\(U\)只是原\(U\)删去已经被本方阻断的线并可能删去不再相关的例外线。真正需要另证的是：原\(U\)尚未阻断的线，怎样在另一方本手之后仍保持空格或固定阻断。

\subsection{3
三类可直接计算的证书}\label{ux4e09ux7c7bux53efux76f4ux63a5ux8ba1ux7b97ux7684ux8bc1ux4e66}

本节只讨论有限\(H\)，不对无限剩余容量取奇偶。记第\(j\)列当前高度为\(h_j\)，余量为\(H-h_j\)。全部证书都必须先满足第2节的共同安全前提。

\subsubsection{\texorpdfstring{3.1 \(A\)类
永久安全或一个目标}{3.1 A类 永久安全或一个目标}}\label{aux7c7b-ux6c38ux4e45ux5b89ux5168ux6216ux4e00ux4e2aux76eeux6807}

满足以下任一项即接受\(A\)类证书。

\begin{itemize}
\tightlist
\item
  \textbf{\(A0\) 永久安全：} \(U\)为空。
\item
  \textbf{\(A1\) 单格封线：}
  有空目标\((j,g)\)，\(1\le g\le H\)，\(g-h_j\)为正偶数，\(j\)列以外总余量为偶数，而且\(U\)的每条线都经过\((j,g)\)。
\item
  \textbf{\(A2\) 下方先胜：}
  目标满足相同的盘内、正偶距、列外偶容量条件。把本方棋放在目标上，会与三个当前已固定的本方棋构成四连；\(U\)的每条线都经过同列某个行号不低于\(g\)的格。
\end{itemize}

\(A2\)只作几何检查：目标四连中的另外三枚必须现在已经是本方棋，不能猜测中间尚空格将来会是什么颜色，也不是允许现在越过空格实际落子。每条线必须被同一个目标保护；这里没有把两条不同线各需一个目标误改成``取得其中一个即可''。

\textbf{引理2 \(A\)类的安全与闭包。}
有\(A\)类证书时，另一方的下一手不能首先获胜。若该手未终局，原列未满，则同列覆盖后仍有\(A\)类证书或本方已胜、盘已满；原列满时，至少有一个合法回应具有同样性质。

\textbf{证明。}
\(A0\)时所有底边、例外线已有本方棋；其余线由支撑下移排除首胜，竖向由第2.3节排除。同列盖或在原列满后任取合法列，都保持此条件。

\(A1\)、\(A2\)时，目标距列高为正偶数。在另一方回合，它最多走到目标下一格，随后本方同列盖合法；所以在本方取得目标之前，目标一直空。\(A2\)同列更高格也因重力一直空。于是\(U\)中的线不能先完成，其他线和竖向仍由共同条件排除。原列未满时，不论它是不是目标列，都同列盖，目标距离要么不变、要么减二。

另一方填满的列不可能是尚待兑现的目标列。填满发生在目标列外，而列外总余量在其回合开始为偶数；只要另一方能落，落前至少两格，落后仍有合法的外部回应。选任何合法的非目标列加本方棋，不会损坏目标、固定棋子或安全条件。

目标未取得时保持原证书；\(A1\)取得时，\(U\)全部被阻断，转\(A0\)；\(A2\)取得时，本方立即四连并停止。任何其他回应若已经提前四连，也立即停止；无胜满盘直接为和棋。故逐步闭包成立。\(\square\)

证书可在每轮重新选择，不必把先前目标保存在策略状态中。旧证书保证另一方本手不能首胜，并保证存在安全后继；实际选择的另一个合格后继再由它自己的证书接管。

\subsubsection{\texorpdfstring{3.2 \(B\)类
零或两列奇余量的覆盖证书}{3.2 B类 零或两列奇余量的覆盖证书}}\label{bux7c7b-ux96f6ux6216ux4e24ux5217ux5947ux4f59ux91cfux7684ux8986ux76d6ux8bc1ux4e66}

定义奇余量列集合

\(R=\{j\)：\(H-h_j\)为奇数\}。

定义保护格集合\(G\)：本方已经占有的全部格，加上所有盘内空格\((j,g)\)，其中\(g-h_j\)为正偶数。这些空格在连续同列覆盖下，只可能保持为空，或由本方取得。

\(B\)类的完整条件为：

\begin{enumerate}
\def\labelenumi{\arabic{enumi}.}
\tightlist
\item
  \(|R|\)为零或二。
\item
  \(U\)中的每条线与\(G\)相交。
\item
  若\(R=\{f,t\}\)，必须对两个有序选择\((f,t)\)、\((t,f)\)分别检查：\(U\)的每条线都与``\(G\)中不在待换入列\(t\)的格，连同该列顶格\((t,H)\)''相交。
\end{enumerate}

第3项不是只检查一条预期路径，而是检查两种首次满列顺序。待换入列的顶格此时一定为空，因为该列余量为正奇数。这个条件不能省略，也不能只检查原\(G\)。

\textbf{引理3 \(B\)类的安全与闭包。}
有\(B\)证书时，另一方下一手不能首胜。原列未满就同列盖，得到\(B\)证书或终局；若另一方填满一列，则\(R\)原来必有两列，在另一奇余量列回应合法，并得到\(|R|=0\)的\(B\)证书或终局。

\textbf{证明。}
\(|R|=0\)时，每个可进入列在另一方回合开始至少剩两格，所以同列覆盖始终合法。\(G\)中的空保护格被填到时一定归本方，未填到则仍空。\(U\)不能被另一方完成；其余线和竖向由第2节处理。

\(|R|=2\)时，完整同列覆盖减去两格，\(R\)不变，\(G\)只会把空保护变为本方固定格而不会丢失保护。完整例外也不增加，所以三项\(B\)条件逐步保持。另一方首次填满的列只能是\(f\in R\)。它填满这一手仍受当前\(G\)保护，不能等本方换列之后才检查其是否已经胜利。

设另一奇余量列为\(t\)，其余量是正奇数，本方在\(t\)落一枚合法。\(f\)余量变零，\(t\)余量由奇减一变偶，其他列原来已偶，故所有余量此后为偶数。

不在\(t\)的原\(G\)保护仍为空或本方；\(f\)已满，其中原来的空\(G\)格已经固定属于本方。\(t\)列的旧空偶距保护可能失效，所以不能继续沿用它们。取而代之，\(t\)顶格\((t,H)\)在这次落子后要么立即属于本方，要么与新列高相差正偶数，进入新的\(G\)。第3项恰保证\(U\)仍有这些有效保护格之一。换列不会引入\(E\)外的新例外；新的\(U\)又只是原\(U\)的子集，因此取得\(|R|=0\)的\(B\)证书。

以后不再需要满列换列。每一枚真实落子后都优先判断首胜或满盘，无须在终局后维护证书。\(\square\)

这是真正的逐步证书闭包，不必记住原来的\(R\)，也不必保存``是否已经换过列''。

\subsubsection{\texorpdfstring{3.3 \(C\)类
两中列全空的等待证书}{3.3 C类 两中列全空的等待证书}}\label{cux7c7b-ux4e24ux4e2dux5217ux5168ux7a7aux7684ux7b49ux5f85ux8bc1ux4e66}

\(C\)类要求：

\begin{itemize}
\tightlist
\item
  四列余量均为奇数。
\item
  第二、第三列全空。
\item
  两外列均为从另一方开始、本方结束的交替偶长列，允许一列为空，但不能都空。
\end{itemize}

例如检查后手证书时，两外列是\((\text{先后})^a\)、\((\text{先后})^d\)，其中\(a\)、\(d\)非负且\(a+d\ge 1\)。检查先手证书时对应\((\text{后先})^a\)、\((\text{后先})^d\)。因为中列为空，四奇余量自动推出\(H\)为奇数；非空外列的偶长又推出\(H\ge 3\)。

\textbf{引理4 \(C\)类的安全与出口。}
另一方在外列落子而未填满时，同列覆盖保持\(C\)类或终局。另一方首次进入中列或填满外列时，有一个合法回应得到\(B\)类或终局。另一方的事件手本身也不能先胜。

\textbf{证明。}
等待中两中列为空，事件手后至少还留一条空中列，故另一方不可能非竖直四连；外列棋子交替，进入中列时该列又只有一枚棋，所以竖向也不可能首胜。外列正常覆盖保持偶长交替，四奇余量不变。

若另一方进入第二列，本方在第三列底部回应；进入第三列则在第二列底部回应。两中列都高一，余量偶数，两外列是余下的奇余量列。底横线含本方棋。唯一潜在例外在另一方中列的第二格，其第二行横线有某条非空外列第二格的本方棋阻断。下降底斜线由\((2,3)\)保护，上升底斜线由\((3,3)\)保护；两格均与中列现高相差二，不在任何待换入外列中。因此\(B\)的原保护和两种满列顺序条件均成立。\(H=3\)时两条底斜线越界，只需底横和例外第二行已经阻断。

若另一方填满第一列，本方在第三列底部回应。\(H\ge 5\)时第一列第四格本来就是本方棋，阻断下降底斜线；上升底斜线由\((3,3)\)保护。奇余量列是第二、第四列，保护不在其中，故两种满列条件均成立。满的另一方顶列不产生盘外潜在例外。\(H=3\)时没有斜四，底横已有本方棋，条件同样成立。第四列填满时，反射地在第二列底部回应。

所有回应都是空中列的底格，合法；其后若未终局即具有\(B\)证书。\(\square\)

这张事件表只是证明合格回应存在，不是统一算法在运行中调用的动作表。

\subsubsection{3.4
为什么先手不会误用后手证书}\label{ux4e3aux4ec0ux4e48ux5148ux624bux4e0dux4f1aux8befux7528ux540eux624bux8bc1ux4e66}

先手回应后，总子数奇，总空格数也奇，所以奇余量列数为奇数。\(B\)要求零或二条，\(C\)要求四条，因此\(B\)、\(C\)在合法先手回应后的局面自动不成立。先手只会被\(A\)类接受。

后手回应后，总空格数为偶数，三类均可直接检查。\(A1\)、\(A2\)列外偶容量进一步推出：先手目标列余量为奇数，目标行与\(H\)奇偶相反；后手目标列余量为偶数，目标行与\(H\)奇偶相同。算法直接检查列外余量，无须写这两套目标行规则，也无须按先后手选择不同的证书库。

\subsection{4
完全相同的确定性回应算法}\label{ux5b8cux5168ux76f8ux540cux7684ux786eux5b9aux6027ux56deux5e94ux7b97ux6cd5}

\subsubsection{4.1
合格回应的定义}\label{ux5408ux683cux56deux5e94ux7684ux5b9aux4e49}

对方刚落以后，先检查是否已有首胜或满盘；若有，停止。否则对每个合法落列\(j\)，暂时放入当前回应一方的一枚棋，得到候选局面。以下任一情形称\(j\)为合格回应：

\begin{enumerate}
\def\labelenumi{\arabic{enumi}.}
\tightlist
\item
  该枚使当前回应一方立即四连。
\item
  没有胜者而直接满盘和棋。
\item
  候选非终局，并通过\(A\)、\(B\)、\(C\)中的任一种证书检查。
\end{enumerate}

计算候选不是真实落子，也不允许在对方已经胜后测试补救。\(A\)、\(B\)、\(C\)的先后检查顺序不影响``至少有一种''的判定。

\subsubsection{4.2
唯一的优先级}\label{ux552fux4e00ux7684ux4f18ux5148ux7ea7}

令\(c\)为上一手列，\(r\)为上一手行。对先手、后手都执行以下同一个顺序。

\textbf{第一步。}
若\(c\)仍可落，并且同列回应合格，或者两个底角\((1,1)\)、\((4,1)\)都已经属于当前回应的一方，就在\(c\)同列覆盖。这里``\(c\)仍可落''适用于两个并列条件。

\textbf{第二步。}
若第一步未选择动作，而存在合格回应：先选立即获胜的回应，多个时取最左列；若没有立即胜，优先选落子前余量\(H-h_j\)为奇数的合格列；仍有并列时取最左列。这里的列高与余量均在对方刚落以后、本方尚未回应的棋盘上计算。

\textbf{第三步。}
若前两步都没有选择动作，在算法可达历史上，\(c\)必未满。若\(r\ge 3\)，就同列盖。若\(r\le 2\)，对每个合法回应计算其后底部线中有多少条含本方棋，取数量最多者；并列时优先同列盖，仍并列时取最左列。

为使回应函数对所有合法非终局输入全定义，第三步先加一项：若\(c\)已满，取最左合法列；否则继续上述行号与评分规则。下文证明从空盘遵照算法的任何路径都不会触发这一补足。它不判断历史是否相容，不读取阶段，只补齐其他输入上的定义。

先手首子是唯一没有上一手的例外，固定占\((1,1)\)；\(H=0\)不落子。此后两方没有不同的回应分支。

\subsubsection{4.3
优先级中的细节}\label{ux4f18ux5148ux7ea7ux4e2dux7684ux7ec6ux8282}

同列合格或双底角条件可以优先于别列立即获胜。不败策略不必抓住每个可赢机会；在尚无证书的偶高先手阶段，这一优先覆盖正用于避免第二次主动换列。

第三步评分是候选之后已经永久阻断的底部线数量，不是估计未来胜率。越界底斜线不计。计算只包括四个合法候选的一手落子、当前四连、例外、目标、余量和保护集合，不包含递归求解或按高度索引的策略图。固定四列且取格视为常数时间时，朴素检查器可枚举\(O(H)\)条线和\(O(H)\)个目标，在\(O(H^2)\)时间内完成一次决策。这是上述取格模型下的算法估计，未证明参考
Lean 程序本身的机器成本。

算法可以每次选择新的证书或目标。正确性依靠证书的逐步闭包，而不是记住最初选中的证书然后照表执行。

\subsection{5
先手全部有限高度的证明}\label{ux5148ux624bux5168ux90e8ux6709ux9650ux9ad8ux5ea6ux7684ux8bc1ux660e}

\subsubsection{5.1
已有证书后不再主动换列}\label{ux5df2ux6709ux8bc1ux4e66ux540eux4e0dux518dux4e3bux52a8ux6362ux5217}

先手只会有\(A\)证书。由引理2，后手下一手不能首胜；若仍需回应，原列未满时同列盖必合格，第一步必覆盖。原列满时至少有一个合格回应，第二步有定义。无论第二步按奇余量和最左规则选哪一个，后继本身都有证书或已经终局。

因此一旦取得\(A\)证书，后面不再进入第三步，也不再主动换列。下面只需证明开局已持证，或无证书阶段仍能安全继续。

\subsubsection{5.2
小高度与奇高度}\label{ux5c0fux9ad8ux5ea6ux4e0eux5947ux9ad8ux5ea6}

\(H=0\)直接结束。\(H=1\)、2、3时，先手首子已经阻断底横线，斜四和竖四不存在，且没有后手例外，所以首子后有\(A0\)。

奇\(H\ge 5\)时，先手首子阻断底横线和上升底斜线，只剩\(D=((1,4),(2,3),(3,2),(4,1))\)。目标\((3,2)\)距空第三列二格，列外余量为\(3H-1\)，偶数。故首子后已有\(A1\)，整个后续只用闭包，主动换列为零。

\subsubsection{5.3
偶高无证书阶段的第一种形状}\label{ux5076ux9ad8ux65e0ux8bc1ux4e66ux9636ux6bb5ux7684ux7b2cux4e00ux79cdux5f62ux72b6}

设偶\(H\ge 4\)。在尚未有证书的先手局面，第一种可能形状为

\[[\text{先}(\text{后先})^a,\varnothing ,\varnothing ,(\text{后先})^d]\]

其中\(a\)、\(d\)非负，\(a\ge 1\)蕴含\(d\ge 1\)。根是\(a=d=0\)。此时只有第一列底角属于先手。下述分类只用于证明，不作为算法状态。

后手每手落后至少有一条空中列，故不能非竖直首胜；外列后手棋相隔，也不能竖四。

\textbf{进入第二列。}
先手在第一列加一枚是一个合格回应见证。得到\([\text{先}(\text{后先})^a\text{先},\text{后},\varnothing ,(\text{后先})^d]\)。唯一潜在例外\((2,2)\)的第二行，\(a=0\)时有\((1,2)=\text{先手}\)，\(a\ge 1\)时有\(d\ge 1\)保证\((4,2)=\text{先手}\)；另一有效斜线已有\((1,1)\)。底横线、上升底斜线也已挡。若\(a=1\)，下降线\(D\)已被刚取得的\((1,4)\)阻断，有\(A0\)；否则\((2,3)\)距第二列现高一为二，列外余量偶，是\(D\)上的\(A1\)目标。回应合法，因为第一列原为小于偶\(H\)的奇高。

\textbf{进入第三列。}
同列盖到\((3,2)\)阻断\(D\)；各非空列由先手封顶，没有例外，有\(A0\)。

\textbf{填满第一列。}
先手在第二列底格回应，目标\((2,3)\)距现高一为二。唯一可能未挡线为\(D\)，其余列偶高或已满，列外余量偶，有\(A1\)。满列顶后手受先手支撑，不新增盘外例外。

以上都只证明至少一个合格回应存在，实际算法可选另一合格回应。只要进入证书，闭包即接管，不必回到原等待。

其余是外列未满的普通手。若同列回应合格，第一步覆盖并进入证书；若另一处已有合格回应，第二步可以提前进入证书。若没有任何合格回应，第三步如何选择，只须核对以下穷尽情形。

\begin{itemize}
\tightlist
\item
  后手在外列落到第三行以上：第三步同列覆盖，保持该形状。
\item
  后手首次落第四列底部：必有\(a=d=0\)。所有合法回应都仍只阻断原来的两条底部线，评分相同，选同列覆盖，得到\(d=1\)。
\item
  后手在第一列第二格落子且\(d\ge 1\)：没有一枚合法回应能新增阻断\(D\)，评分相同，选同列覆盖，得到\(a=1\)。
\item
  剩下唯一低行情形是根上后手落第一列第二格。只有回应第四列底格能把三条底部线全挡住，因此若此前仍无合格回应，第三步唯一最高分选择第四列，占齐底角，进入下一种形状。\(H=4\)时可能更早由第二步选取某个合格出口，此时直接进入证书，不要求实际仍走第四列。
\end{itemize}

所以只要仍无证书，该形状在普通覆盖下保持，或仅经一次主动换列转入第二种形状；其他主动回应都直接结束无证书阶段。

\subsubsection{5.4
偶高无证书阶段的第二种形状}\label{ux5076ux9ad8ux65e0ux8bc1ux4e66ux9636ux6bb5ux7684ux7b2cux4e8cux79cdux5f62ux72b6}

第二种形状为

\[[\text{先后}(\text{后先})^a,(\text{后先})^b,(\text{后先})^e,\text{先}(\text{后先})^d]\]

这里把第三列参数写为\(e\)，避免与上一手列\(c\)混淆。前三列偶高，第四列奇高，两个底角都属于先手。后手只可能填满第四列。

底部三线已由两个底角阻断。唯一已有或潜在例外是\((1,3)\)：第三行由始终为空或先手的\((4,3)\)保护，\(U_1=((1,3),(2,4),(3,5),(4,6))\)由始终为空或先手的\((2,4)\)保护。第一列顶端后手串最多长到两枚即被盖，其他后手棋相隔。其余线由支撑下移排除首胜。

原列未满时，第一步的双底角条件保证覆盖，即使另一列能提前胜或取得证书，也不会再次主动换列。同列覆盖保持形状，或提前进入\(A\)类闭包。

每轮消耗两格，若此前未终局或进入证书，后手最终填满第四列。若该手已经满盘，直接结束；否则下表给出至少一个合格回应。所有条件互斥、穷尽。

{\def\LTcaptype{none} % do not increment counter
\begin{longtable}[]{@{}
  >{\raggedright\arraybackslash}p{(\linewidth - 2\tabcolsep) * \real{0.5000}}
  >{\raggedright\arraybackslash}p{(\linewidth - 2\tabcolsep) * \real{0.5000}}@{}}
\toprule\noalign{}
\begin{minipage}[b]{\linewidth}\raggedright
条件
\end{minipage} & \begin{minipage}[b]{\linewidth}\raggedright
合格回应的存在见证
\end{minipage} \\
\midrule\noalign{}
\endhead
\bottomrule\noalign{}
\endlastfoot
\(H=4\) &
任一合法列；\(U_1\)越界，第三行已有\((4,3)=\text{先手}\)，得到\(A0\)。 \\
\(H\ge 6\)，\(a=0\) & 占\((1,3)\)，消除潜在例外，得到\(A0\)。 \\
\(H\ge 6\)，\(a\ge 1\)，\(b\ge 2\) &
任一合法列；\(U_1\)已有\((2,4)=\text{先手}\)，得到\(A0\)。 \\
\(H\ge 6\)，\(a\ge 1\)，\(b\le 1\)，\(e=2\) &
占\((3,5)\)，封\(U_1\)，得到\(A0\)。 \\
\(H\ge 6\)，\(a\ge 1\)，\(b\le 1\)，\(e\le 1\) &
在第三列加一枚；\((3,5)\)为\(A1\)目标。 \\
\(H\ge 6\)，\(a\ge 1\)，\(b=1\)，\(e\ge 3\) &
占\((2,3)\)，在\(D\)上立即获胜。 \\
\(H\ge 6\)，\(a\ge 1\)，\(b=0\)，\(e\ge 3\) &
占\((2,1)\)；\((2,3)\)为\(A2\)制胜目标。 \\
\end{longtable}
}

\(a\ge 1\)时例外\((1,3)\)实际存在，其横线已被满第四列的\((4,3)\)阻断，只可能剩\(U_1\)。\(e\le 1\)时，先手在第三列加至\(2e+1\)，目标\((3,5)\)距新列高\(4-2e\)，为正偶；其他列偶高或已满，列外余量偶。\(e\ge 3\)时，\(D\)的\((1,4)\)、\((3,2)\)、\((4,1)\)都是先手，故第二列高二可以立即赢；第二列空则先占底格，再在第三格先赢，\(U_1\)要求的同列第四格还在其上方。

指定落点不超过第五行，相关分支\(H\ge 6\)保证盘内。此表是第二步非空的证明，不是算法的七条执行分支；实际算法仍只用相同的证书、奇余量及最左优先级。

\subsubsection{5.5
先手归纳结论}\label{ux5148ux624bux5f52ux7eb3ux7ed3ux8bba}

小高度和奇高度从首子即持证。偶高的无证书局面只可能属于上面两种形状；每个后手手先被证明不能首胜，随后算法要么安全保持形状，要么取得证书，原列满时必有合格回应。每枚棋消耗一格，故不存在无限等待。先手算法在全部有限高度不败。

无证书第一形状只可能在``进入证书''或``占齐底角转第二形状''时主动换列一次；第二形状以及持证阶段都不再主动换列。因此整局主动至多一次。奇高度及\(H\le 3\)从开局持证，主动次数为零。

\subsection{6
后手全部有限高度的证明}\label{ux540eux624bux5168ux90e8ux6709ux9650ux9ad8ux5ea6ux7684ux8bc1ux660e}

\subsubsection{6.1
开局总有合格回应}\label{ux5f00ux5c40ux603bux6709ux5408ux683cux56deux5e94}

以下分类只为证明统一检查器能找到合格回应，不是算法里的开局动作表。

\(H=1\)时先手首列已满，后手在任一其他列落底。底横线有后手阻断，没有有效例外，有\(A0\)；另外两空列为奇余量，亦满足\(B\)。

奇\(H\ge 3\)、先手首下外列，后手同列盖后，两外列为``先后''偶长交替、两中列空、四余量奇，有\(C\)。同列合格优先，算法确实选择覆盖。

奇\(H\ge 3\)、先手首下第二列，后手占第一列底格是\(B\)见证：\(R=\{3,4\}\)，底横和上升底斜线已挡；\(D\)由第二列第三格保护，该格不在\(R\)中。潜在例外\((2,2)\)的第二行由第三、第四列第二格共同保护；若将来换入其中一列，另一列的第二格仍留在\(G\)里，满足两种有序满列条件。第一列开局棋不产生先手例外。首下第三列由反射得到同样见证。

正偶\(H\)时，可在相邻伙伴列回应，配对为\{1,2\}、\{3,4\}。\(R\)为开局落过子的两列，其余两空列相邻且余量偶。只须核对首下第一、第二列；另两种反射。

\begin{itemize}
\tightlist
\item
  先1、后2：底横有后手；两条底斜线可分别用第四列第四格、第三列第二格保护。唯一潜在例外\((1,2)\)的第二行用第三列第二格保护，可能存在的上升线经过第三列第四格。保护均在两个空列，不在待换入的\(R\)列。
\item
  先2、后1：底横与上升底斜线已经含后手；\(D\)用第三列第二格保护。潜在例外\((2,2)\)的第二行也用第三或第四列第二格保护；其另一有效斜线已有后手(1,1)。
\end{itemize}

这些空列的偶数行均与现高零相差正偶数，两种满列顺序条件成立。\(H=2\)时只检查横线，所有斜线越界。因此每个后手开局都存在合格回应。

\subsubsection{6.2
从证书到同一算法}\label{ux4eceux8bc1ux4e66ux5230ux540cux4e00ux7b97ux6cd5}

\(A\)、\(B\)局面原列未满时都有合格覆盖，算法第一步必选覆盖。原列满时都有合格回应，第二步有定义。\(C\)局面外列未满时有合格覆盖；首次中列或满外列事件有引理4给出的\(B\)回应，第二步也有定义。

在\(C\)局面，两外列底格属于先手，不可能是后手的两个底角。因此第一步的双底角条款不会强迫后手在中列事件错误覆盖。\(A\)、\(B\)阶段原本就有合格覆盖，双底角条件也不会改变动作。后手初次回应前还没有任何后手底角。

于是后手从首次回应起始终持有证书，从不使用第三步的低行评分或高行兜底。每个先手手先由当前证书排除首胜，之后再执行闭包或事件出口。后手算法在所有有限高度不败。

\subsubsection{6.3
为什么可以同时保留最多两次总换列}\label{ux4e3aux4ec0ux4e48ux53efux4ee5ux540cux65f6ux4fddux7559ux6700ux591aux4e24ux6b21ux603bux6362ux5217}

还需证明引入\(A\)类和任意合格出口后，计数不会暗中增加。后手每次回应后总余量为偶数。

\textbf{首次进入\(A\)或\(B\)时，奇余量列至多两条。}
正偶开局若同列盖，四列余量都偶，但底横没有后手棋，也没有能以正偶距保护底格的空目标；\(A\)、\(B\)均不能通过，\(C\)也不适用，所以合格回应必须异列，得到两条奇余量列。奇高中列开局若同列盖，则四余量都奇、底横仍无后手棋、两中列并非全空，\(A\)、\(B\)、\(C\)均不能通过，所以合格回应也必异列，得到两奇余量。

奇高外列开局同列盖进入\(C\)。\(C\)阶段普通外列覆盖仍留四奇余量，且底横没有后手棋；\(A0\)不能成立，\(A\)单目标的列外三条奇余量之和为奇数，也不能成立，\(B\)要求的条数不符，所以不会误转进四奇余量的\(A\)或\(B\)。中列事件后若继续同列盖，仍有四奇余量、无后手底格且\(C\)不成立，因此该覆盖不合格；空的另一中列也排除这枚覆盖直接获胜。合格异列回应使两条奇余量一起变偶，得到两条。外列满事件也先把一条余量奇变零，再由异列回应消去另一条，仍剩两条。\(H=1\)的首次回应同样留下两空列奇余量。

\textbf{进入\(A\)或\(B\)之后，奇余量不会增加。}
原列未满时算法始终覆盖，每列余量减二，\(R\)不变。若先手填满原列，\(R\)此前必为两条，另一奇余量列唯一。

\(B\)证书保证在这另一列回应合格。\(A0\)允许任一合法列，也包括它。\(A1\)、\(A2\)的目标列余量必为偶数，因为后手总余量偶且列外余量偶；所以另一奇余量列不可能是目标列，在那里回应正是合法安全的目标列外回应。若没有立即后手胜，第二步的奇余量优先就必选此唯一奇余量列，使所有余量变偶。此后先手不能再填满任何可进入列，永久不再换列。若优先选择某个立即胜，游戏已经结束，更不需要后续计数。

所以后手的主动换列只可能是开局异列，或开局同列进入\(C\)以后首次中列事件的异列，两者不在同一路径重复发生。最多一次。开局或\(C\)首事件至多贡献一次总换列，随\(\text{后}A/B\)阶段至多一次被动换列，故整局总换列至多两次。\(H=1\)开局列已满，首次换列是被动而不是主动。

这段证明需要``奇余量优先''。若任取合格满列回应，安全性仍可能成立，但``换列后全部余量偶、之后不再换列''的次数论证不能原样沿用。

\subsection{7
总定理与最优主动次数}\label{ux603bux5b9aux7406ux4e0eux6700ux4f18ux4e3bux52a8ux6b21ux6570}

\textbf{定理5 统一回应。}
对四列及任意有限\(H\ge 0\)，\(H=0\)直接和棋；\(H>0\)时先手首下第一列，此后先手和后手均执行第4节同一个回应函数。各方分别面对任意合法对手续法都不败。先手、后手整局主动换列至多一次；先手总换列至多四次，后手至多两次。

\textbf{证明。}
不败由第5、6节。双方各自用不败策略时，不能有任何一方首先获胜，有限容量迫使无胜满盘，故空盘值为和棋。主动次数已分别证明。每次被动换列对应对方刚填满的一列，同一列不能重复填满；第四列也满时全盘终局，没有第四次回应，因此任一方被动至多三次，先手总数至多四。后手较强的两次界由第6.3节给出。\(\square\)

\textbf{定理6 主动预算最优。}
在四列有限盘上，第4节算法达到以下最优最坏情形主动预算：

{\def\LTcaptype{none} % do not increment counter
\begin{longtable}[]{@{}lrr@{}}
\toprule\noalign{}
高度 & 先手 & 后手 \\
\midrule\noalign{}
\endhead
\bottomrule\noalign{}
\endlastfoot
\(H=0\)、1 & 0 & 0 \\
\(H=2\)、3 & 0 & 1 \\
奇\(H\ge 5\) & 0 & 1 \\
偶\(H\ge 4\) & 1 & 1 \\
\end{longtable}
}

这里的一次是主动预算的最优值，不是总换列次数、公式长度或证书数量的最优值。下节连接配套数学论文的完整下界；上界已由本算法证明。

\subsection{8
主动预算的最优性接口}\label{ux4e3bux52a8ux9884ux7b97ux7684ux6700ux4f18ux6027ux63a5ux53e3}

配套数学论文 {[}4, §7{]}
证明，后手在所有有限\(H\ge2\)时不能以零主动预算保证不败；先手在所有偶\(H\ge4\)时同样不能。后手下界由依次抢占四个底格给出。先手下界分别处理边列和中列首手，允许满列以后任意依赖历史的被动选择；其中中列陷阱的第一列准备高度必须恰为四。该下界不是对某一贪心换列规则的反例。

结合本文算法的上界，在四列有限盘上，先手最优主动预算为\(1\)当且仅当\(H\ge4\)且\(H\)偶；后手最优值为\(1\)当且仅当\(H\ge2\)；其他情形均为零。详尽的参数化反制保留在数学论文，本文不重复。Lean
的完整历史语义与预算递归等价，因此同一结论覆盖任意历史策略，不限于无额外状态的实现。

\subsection{9
证书条件为什么不能随意删去}\label{ux8bc1ux4e66ux6761ux4ef6ux4e3aux4ec0ux4e48ux4e0dux80fdux968fux610fux5220ux53bb}

\subsubsection{9.1
只检查原保护集合并不安全}\label{ux53eaux68c0ux67e5ux539fux4fddux62a4ux96c6ux5408ux5e76ux4e0dux5b89ux5168}

取\(H=3\)、轮到先手，局面为

\[[\text{先先},\text{后先后},\text{后先后},\varnothing ]\]

它可由合法棋谱``先1，后2，先1，后3，先2，后2，先3，后3''达到，之前无人四连。对于后手的检查，奇余量列是第一、第四列，唯一未挡线为第二行。第四列第二格与现高零相差二，属于原\(G\)。因此若删掉\(B\)的两种满列后条件，只看原保护，会错误接受。

但接着先手填第一列顶格，后手转第四列底格，先手再占第四列第二格，就横四首胜。原偶距保护在换入第四列时失效。正确的\(B\)条件要求第二行与``不在第四列的\(G\)，或第四列顶格''相交，恰好拒绝该局面。

\subsubsection{9.2
安全选择不保证棋谱相同}\label{ux5b89ux5168ux9009ux62e9ux4e0dux4fddux8bc1ux68cbux8c31ux76f8ux540c}

取 \(H=3\)
的合法前缀``先1，后2，先2，后2''。第二列刚满。若在合格回应中取最左列，会选择第一列第二格；本文规则优先奇余量列，选择第三列底格。两手都留下永久安全证书，但并非同手。

因此，安全性不能代替两个选择器的逐手等价证明。本文证明所定义规则本身的安全和次数，不声称它与配套数学构造允许的每一种选择完全同手。

\subsection{10
交替段与目标扫描}\label{ux4ea4ux66ffux6bb5ux4e0eux76eeux6807ux626bux63cf}

\subsubsection{10.1
候选棋盘至多九段}\label{ux5019ux9009ux68cbux76d8ux81f3ux591aux4e5dux6bb5}

最大交替段是列中相邻棋子颜色交替的极大连续区间，同色相邻的两格在其间切开。总段数\(K\)等于非空列数加同色邻接对数；它不是同色棋串数。

\textbf{命题7 候选段数界。}
设宽度\(w>0\)，棋盘从空盘开始合法交替落子。某受控方的实际主动次数至今为\(A\le a\)，被动次数为\(B\)。在对方刚落、尚未终局、受控方即将回应时，任意合法候选回应所得棋盘满足

\[K'\le 2w-1+2a.\]

\textbf{证明。}
受控方的同色邻接对至多\(A+B\)：同列覆盖不产生它，免费的首子也不产生它。对方每个同色邻接对，可记在其下方那枚棋子刚落后受控方没有盖住的回应上。该列当时尚未满，否则上方不能再有棋子，所以这是主动回应；不同邻接对对应不同的下格，因而对应不同回应。这也计入对方刚形成的新邻接对。因此对方的同色邻接对至多\(A\)，回应前\(K\le w+2A+B\)。

每次被动换列封掉不同的一列，且当前仍要回应，所以\(B\le w-1\)。若\(B\le w-2\)，候选落子至多增加一段，于是\(K'\le w+2a+w-2+1\)。若\(B=w-1\)，至少已有\(w-1\)列满；当前唯一合法列就是对方刚落的那列。任何候选回应都盖在异色棋上，不增加段，同样有\(K'\le2w-1+2a\)。证明不要求尚未选择的假想回应也遵守主动预算。\(\square\)

取\(w=4\)、\(a=1\)，统一界为九。这个界在所用规则中可以达到。高度六、先手受控时，以下列号序列全程合法且没有提前终局：

\texttt{1,2,1,1,1,2,2,2,2,2,1,3,3,3,3,4,4,1,3,3,4}。

记受控棋为\(P\)、对手为\(Q\)，最终四列是\([PPQPPQ,QQPQPQ,QPQPPQ,QPP]\)，共有九个最大交替段，主动一次、被动三次。它既是实际回应后的状态，也是该步的合法候选。随附坐标脚本重放了所有前缀。该实例证明本段数界不能统一降至八，不证明总换列四次对所有不败策略都必要。

\subsubsection{10.2
制胜目标的有序区间扫描}\label{ux5236ux80dcux76eeux6807ux7684ux6709ux5e8fux533aux95f4ux626bux63cf}

设目标列\(j\)现高\(h_j\)，候选目标行\(g\)须满足\(g\ge h_j+2\)及\(g\equiv h_j\pmod2\)，上界\(u\)由棋盘高度和未决线路在该列的交点给出。目标至少比当前列顶高二，故竖直线不可能已有另外三枚固定己方棋；只须考虑非竖直斜率\(\delta\in\{-1,0,1\}\)。

其他列\(i\)的一个交替段\([l,r]\)内，己方棋占一个确定奇偶类。目标线要求该列的\(g+\delta(i-j)\)为己方棋，因此这一段对\(g\)提供区间

\[[l-\delta(i-j),\ r-\delta(i-j)].\]

过滤与目标奇偶不相容的段，得到三个有序、不相交的区间序列。用三个指针与\([h_j+2,u]\)求交：取当前左端最大值\(L\)、右端最小值\(R\)，把\(L\)上调到所需奇偶的首个整数；若不超过\(R\)，找到目标，否则推进右端最小的区间。被推进的区间与任何更高的后续区间都不可能产生遗漏解。每个指针只向前，故固定\(j,\delta\)的扫描为\(O(K)\)，四个列与三个斜率仅增加常数因子。

独立逐格枚举与新子程序进行了80,020次对照，全部一致，另检查高度\(10^{12}\)、仅三段的压缩输入。这里只交付了
A2 目标扫描子程序及其证明；完整压缩表示精化、与 \texttt{respond}
逐手一致及整个程序的复杂度尚未形式化。在线沿已证策略维护常数个段，与从原始\(H\)行棋盘重新读入是不同任务：前者可用常数个高度字，后者至少要读取输入。每个高度字需要\(\Theta(\log H)\)位，不能把字操作常数时间直接称为常数位复杂度。

\subsection{11
形式化范围与复现}\label{ux5f62ux5f0fux5316ux8303ux56f4ux4e0eux590dux73b0}

\subsubsection{11.1
游戏模型和历史费用}\label{ux6e38ux620fux6a21ux578bux548cux5386ux53f2ux8d39ux7528}

随附工程采用独立的四列有限棋盘模型，以相对颜色\(P\)、\(Q\)描述受控方与对手，严格区分先手免费首子和其后的回应。预算递归在对方每个合法选择之后检查对方首胜、满盘和实际回应；不能在对方已赢后继续补救，也不能漏收己方获胜回应的费用。

\texttt{HistoryBudget.lean} 定义合法历史前缀及其费用语义，证明
\texttt{budgetSafe\_iff\_historyBudget}。\texttt{historyController}
可以保存完整落列序列，策略函数不受有限记忆或可计算性限制。由此，精确预算不仅属于某种有限控制状态，也适用于任意完整历史策略。非相容历史可以任意补成全定义，安全和费用只要求在真实可达历史上成立。

\subsubsection{11.2
根定理与未覆盖范围}\label{ux6839ux5b9aux7406ux4e0eux672aux8986ux76d6ux8303ux56f4}

核心策略与预算声明在 \texttt{UniformFour}
命名空间下；完整名称和类型以源码为准。

{\def\LTcaptype{none} % do not increment counter
\begin{longtable}[]{@{}
  >{\raggedright\arraybackslash}p{(\linewidth - 4\tabcolsep) * \real{0.3333}}
  >{\raggedright\arraybackslash}p{(\linewidth - 4\tabcolsep) * \real{0.3333}}
  >{\raggedright\arraybackslash}p{(\linewidth - 4\tabcolsep) * \real{0.3333}}@{}}
\toprule\noalign{}
\begin{minipage}[b]{\linewidth}\raggedright
结论
\end{minipage} & \begin{minipage}[b]{\linewidth}\raggedright
主要源码入口
\end{minipage} & \begin{minipage}[b]{\linewidth}\raggedright
形式化范围
\end{minipage} \\
\midrule\noalign{}
\endhead
\bottomrule\noalign{}
\endlastfoot
四列有限高度、先手和后手的同一规则不败 &
\texttt{UniformFour/FirstRoot.lean}、\texttt{SecondRoot.lean} &
参数化证明已验收 \\
最优主动预算的上界、下界 &
\texttt{FirstExact.lean}、\texttt{SecondExact.lean} &
参数化证明已验收 \\
完整历史语义及精确预算 &
\texttt{HistoryBudget.lean}、\texttt{HistoryRoots.lean} &
双向等价和根定理已验收 \\
释放容量、两奇列转移、区间奇偶 & \texttt{ReserveArithmetic.lean} &
算术引理已验收 \\
九段界 & \texttt{SegmentBoundArithmetic.lean} &
只验收最后的计数算术，事件单射等前提显式保留 \\
完整多目标通用引理、总换列上界 & 配套数学论文 & 纸面证明 \\
小于四列及真正无限高度的和棋结论 & \texttt{proof/Narrow} &
独立模型已验收，不属于本有限算法 \\
\texttt{Covering} 与 \texttt{Narrow} 两种棋盘模型的等价 & 两工程之间 &
尚未形式化 \\
完整压缩实现及成本 & 第10节 & 尚未完成程序精化 \\
\end{longtable}
}

两个主要最优值声明为
\texttt{UniformFour.U1.first\_exact\_history\_budget} 和
\texttt{UniformFour.U1.second\_exact\_history\_budget}。形式化得到的数值正是第8节的最优值；不以有限枚举替代参数化证明。

\subsubsection{11.3
验收与辅助实验}\label{ux9a8cux6536ux4e0eux8f85ux52a9ux5b9eux9a8c}

工程使用 Lean 4.34.0-rc2 {[}3{]}，仅依赖随 Lean 提供的库，不使用 Mathlib
或大型枚举证书。55/55模块编译成功，14个重点声明的公理闭包仅含
\texttt{propext}、\texttt{Classical.choice}、\texttt{Quot.sound}。对总入口
\texttt{Audit} 的全部传递导入运行
\texttt{lean\ -t\ 0}，内核重放退出码为零；源文件和构建产物的 SHA-256
在重放前后均与编译回执一致。普通警告不影响这些成功条件。

辅助证据另行记录：本次独立坐标程序检查先手偶数高度4至200的全部四首列，共396个实例，并检查后手高度2至200的199个下界实例。附件提供的参考规则检查器在高度0至12、双方角色下重放83,429个状态、238,014条对手边；这次重用检查器并加入段数统计，不把它称为新写的独立实现。Lean
实际回应又与参考函数作了70个边界实例对照。上述有限结果用于发现实现和边界问题，不承担全高度数学证明。

在仓库根目录执行：

\begin{Shaded}
\begin{Highlighting}[]
\ExtensionTok{python}\NormalTok{ tools/verify\_proof.py}
\ExtensionTok{python}\NormalTok{ tools/verify\_proof.py }\AttributeTok{{-}{-}build} \AttributeTok{{-}{-}replay}
\end{Highlighting}
\end{Shaded}

第一条只核对已发布源码的哈希及验收回执，不运行
Lean。第二条重新构建并重放 \texttt{proof/Covering}
的55个模块。须先安装其 \texttt{lean-toolchain}
指定的版本；可通过环境变量 \texttt{LEAN\_BIN} 提供 Lean 路径。每个 Lean
进程使用运行时读取的全部逻辑核。

\texttt{proof/Narrow}
另行提供一至三列有限高度和一至四列无限高度的和棋定理。其四个数学模块及审计入口已重新编译，六项重点公理闭包符合相同白名单，信任级别零的导入重放通过。该工程需要固定版本的
Mathlib；准备依赖后，执行
\texttt{python\ tools/verify\_proof.py\ -\/-supporting\ -\/-build\ -\/-replay}。两个棋盘模型之间的等价尚未形式化。完整步骤、源码清单与验收回执见
\texttt{docs/verification.md} 和 \texttt{audit/}。

\subsection{12
实现边界与后续工作}\label{ux5b9eux73b0ux8fb9ux754cux4e0eux540eux7eedux5de5ux4f5c}

本次交付的是已验收的有限棋盘参考算法、无损交替段的计数证明和 A2
线性扫描子程序。没有把完整压缩引擎写成已交付软件，也没有证明常数大小证书库或策略表达式已经最小。后续应先证明压缩表示与逐格棋盘的取格、落子、线路和证书检查一致，再证明整个选择器返回同一列，最后核算机器成本。沿特定策略保持常数段数与任意棋盘的导入、查询是不同输入问题。

算法顺序的调整也必须重新核验：立即获胜优先、任意合法满列转移或者删去一项出口条件，都可能破坏已证的费用或安全不变量。第二篇的实现优化不会改变第一篇作为独立数学证明的地位。

\subsection{参考文献}\label{ux53c2ux8003ux6587ux732e}

{[}1{]} L. Victor Allis. \emph{A Knowledge-based Approach of
Connect-Four: The Game is Solved: White Wins}. Master's thesis, Vrije
Universiteit Amsterdam, Report IR-163, 1988. 特别参见 §§2.2, 6.1,
7.1。\href{https://tromp.github.io/c4/connect4_thesis.pdf}{原文}。

{[}2{]} Yoshiaki Yamaguchi, Kazunori Yamaguchi, Tetsuro Tanaka, Tomoyuki
Kaneko. Infinite Connect-Four Is Solved: Draw. In \emph{Advances in
Computer Games}, LNCS 7168, 2012, pp.~208--219.
\href{https://doi.org/10.1007/978-3-642-31866-5_18}{doi:10.1007/978-3-642-31866-5\_18}。

{[}3{]} Leonardo de Moura, Sebastian Ullrich. The Lean 4 Theorem Prover
and Programming Language. In \emph{Automated Deduction -- CADE 28}, LNCS
12699, 2021, pp.~625--635.
\href{https://doi.org/10.1007/978-3-030-79876-5_37}{doi:10.1007/978-3-030-79876-5\_37}。本文所用的具体工具链版本另见验收记录。

{[}4{]} GPT.
\emph{一至四列重力四子棋的和棋策略：覆盖策略、最优主动换列次数与 Lean
验证}. 配套数学论文，2026年11月8日。
\href{https://github.com/wugang6812/narrow-connect-four/blob/main/papers/paper-zh.pdf}{PDF}.

{[}5{]} GPT. \emph{UniformFour: Lean definitions and proofs}.
\href{https://github.com/wugang6812/narrow-connect-four/tree/main/proof/Covering}{Source}.
2026.

{[}6{]} GPT. \emph{Interval scanning and coordinate checks}.
\href{https://github.com/wugang6812/narrow-connect-four/tree/main/experiments}{Source}.
2026.

\end{document}
